\chapter{有向グラフ}
    \section{諸定理}
        \subsection{\texorpdfstring{$G_1=(\nu,\varepsilon_1),G_2=(\nu,\varepsilon_2)$}{G1=(ν,ε1),G2=(ν,ε2)}に対して\texorpdfstring{${G_1}^{|\nu-1|} = {G_2}^{|\nu-1|}$}{G1\string^|ν-1| = G2\string^|ν-1|}であるとき、\texorpdfstring{$G_1$}{G1}が全域木を持つ\texorpdfstring{$\iff\;G_2$}{ ⇔ G2}が〃}
            \begin{proof}
                \begin{align*}
                    G_1\textrm{が全域木を持つ} &\iff ^\exists i\in\nu\st^\forall j\in\nu,\;^\exists(i,j)\in{\varepsilon_1}^{|\nu-1|}\quad(\because\textrm{\cite{マルチエージェントシステムの制御}定理2.17}) \\
                    &\iff ^\exists i\in\nu\st^\forall j\in\nu,\;^\exists(i,j)\in{\varepsilon_2}^{|\nu-1|}\quad(\because\textrm{仮定より}{\varepsilon_2}^{|\nu-1|}={\varepsilon_1}^{|\nu-1|}) \\
                    &\iff G_2\textrm{が全域木を持つ}
                \end{align*}
            \end{proof}
        \subsection{非閉路的ならばpruningのアルゴリズムはデッドロックしない}
            前提として、対象としているグラフD内の任意の点はスタート点から到達可能であるとする。Dが非閉路的ならpruningのアルゴリズムはデッドロック(ある段階で未決定点ばかりになって先に進まなくなること)しないことを示す。
            \begin{proof}
                \quad\par
                対偶「デッドロックするなら閉路が存在する」を示す。デッドロックが起こって$m$個の未知点$u_1,u_2,\dotsc,u_m$が残ったとする。$u_1$に進入する未知点が$\{u_2,\dotsc,u_m\}$の中に少なくとも1つある(そうでなければ$u_1$はもはや未知ではない)。必要ならば$u_2$以降の番号を振り直して、$u_1$に進入する未知点の１つを$u_2$とする。するとやはり$u_2$に進入する未知点が1つ以上ある。もしそれが$u_1$であれば閉路$(u_1,u_2,u_1)$が発生して証明は終了する。そうでないとき必要ならば$u_3$以降の番号を振り直して、$u_2$に進入する未知点の１つを$u_3$とする。するとやはり$u_3$に進入する未知点が1つ以上ある。もしそれが$\{u_1,u_2\}$の中にあれば閉路が発生して証明は終了する。そうでないとき必要ならば$u_4$以降の番号を振り直して、$u_3$に進入する未知点の１つを$u_4$とする...。これを$m-1$回繰り返すと、途中で閉路が発生して証明が終了するか、或いは$u_{m-1}$に進入する未知点が$u_m$であるというところまでくる。やはり$u_m$に進入する未知点が1つ以上あるのだが、それは$u_1,u_2,\dotsc,u_{m-1}$の中のどれかである。もう他に選択肢がない。ここまでの議論で道$(u_{m-1},u_{m-2},\dotsc,u_2,u_1)$が存在するから、$u_m$に進入する未知点がどれであろうと必ず閉路が発生する。
            \end{proof}
        \subsection{スタートからゴールへの全ての道の連鎖重みの総和は1以下}
            \begin{shadebox}
                有向グラフ$G(V,E)$を考える。ある点集合$S = \{s_1,\dotsc,s_{N_s}\}\in V$と、ある一つの頂点$g\in V$を選ぶ。
                $S$と$g$は次の条件を満たすものとする。
                \begin{itemize}
                    \item $S$のある1点から$S$の1点(自分自身も含む)へ至る道は存在しない
                    \item $g$を含む閉路が存在したとしても、その閉路上に$S$の点は存在しない
                \end{itemize}
                このとき、$S$の点から$g$へ至る経路全ての連鎖重みを足し上げると、結果は1以下である。
            \end{shadebox}
            「連鎖重み」というのは、辺を通過するたびにその辺の重みを掛け算していくという意味である。
            $G$に閉路がなければ、例えば$s_1$から$g$へ至るパスは有限通りしかない。
            しかし$g$を含む閉路が存在すると、一度$g$へ到着してから閉路を何度か回って$g$に落ち着く場合があるので、パターンは無限通りある。
            しかし閉路を通る毎に重みが掛け算されてどんどん小さくなっていくので、総和が収束しそうである。
            上の定理の主張は、閉路があっても総和が収束し、しかも1以下になるということである。
            \begin{proof}
                \quad\par
                まず重要な要素を網羅する具体例を挙げ、それを解析して、証明すべき核心要素を炙り出そう。
                どうやら閉路が重要そうだ。
                \begin{figure}[H]
                    \centering
                    \includegraphics[keepaspectratio, scale=0.6]
                    {\currfiledir/imgs/SumOfChainedWeightPath/G_orig.eps}
                    \caption{$G(V,E)$}
                \end{figure}
                上のグラフは閉路を含んでいる。$g\rightarrow 8\rightarrow 9\rightarrow g$はまぁなんとかなりそうだ。
                しかし$g\rightarrow 7\rightarrow 6\rightarrow 5\rightarrow g$では5において4経由で$s_1$からの入り込みがあってヤバい感じがする。
                \par
                今求めたい総和というのは結局次のものと同じである。
                まず$s_1,\dotsc,s_3$に立ったプレーヤーが初めLIFE1を持っており、そこから出発して$g$に向かうとして、辺を通る度にプレーヤーのLIFEに重みを掛け算していき、$g$に落ち着いた時点でのプレーヤーのLIFEを評価する。
                プレーヤーの可能な道の通り方全てについて評価点を足し算した値が今求めたい総和である。
                \par
                このグラフを次のように5段階で変形していく。各段階で何をやっているかは下で説明する。
                \begin{figure}[H]
                    \begin{minipage}{0.33\hsize}
                        \centering
                        \includegraphics[keepaspectratio, scale=0.6]
                        {\currfiledir/imgs/SumOfChainedWeightPath/step1.eps}
                        \caption{step1}
                    \end{minipage}
                    \begin{minipage}{0.33\hsize}
                        \centering
                        \includegraphics[keepaspectratio, scale=0.6]
                        {\currfiledir/imgs/SumOfChainedWeightPath/step2.eps}
                        \caption{step2}
                    \end{minipage}
                    \begin{minipage}{0.33\hsize}
                        \centering
                        \includegraphics[keepaspectratio, scale=0.6]
                        {\currfiledir/imgs/SumOfChainedWeightPath/step3.eps}
                        \caption{step3}
                    \end{minipage}
                \end{figure}
                \begin{figure}[H]
                    \begin{minipage}{0.33\hsize}
                        \centering
                        \includegraphics[keepaspectratio, scale=0.6]
                        {\currfiledir/imgs/SumOfChainedWeightPath/step4.eps}
                        \caption{step4}
                    \end{minipage}
                    \begin{minipage}{0.33\hsize}
                        \centering
                        \includegraphics[keepaspectratio, scale=0.6]
                        {\currfiledir/imgs/SumOfChainedWeightPath/step5.eps}
                        \caption{step5}
                    \end{minipage}
                \end{figure}
                \noindent
                \textbf{Step1}
                \par
                $s_1$から$g$へ至るパスを考えると、4へ着いた時点でのLIFEに、4から$g$までのパスの重みを掛け算することになる。
                だから最初から$s_1$を4に移動させておけばいい。ここでは新しいノード$s_4$で4を置き換えている。
                \par
                $s_2$から$g$へのパスと$s_3$から$g$へのパスのパターンをもし全部列挙したなら、1に着いてから先は両者で共通である。
                今扱っている問題では総和を考えているから、グラフ上で$s_2$と$s_3$を統合しても計算結果は変わらない。
                ここでは新しいノード$s_5$を作っている。\\
                \noindent
                \textbf{Step2}
                \par
                $s_4$から$g$へのパスを考えると、重み1/2の辺を通る場合と1/4の辺を通る場合を別々に考えることができる。
                しかも$s_4$を一度離れるともう戻ってこない。
                だからグラフ上で$s_4$からの行き先の数だけ$s_4$のコピーを作っても計算結果は変わらない。
                ここでは新しく$s_6,s_7$を作っている。\\
                \noindent
                \textbf{Step3}
                \par
                ここが一番ややこしい。$s_7$から$g$へのパスを考えると、$s_5$を経由して$g$に着いてそこでお終いというパターンと、その後で閉路に入って何度か回って$g$に落ち着くパターンがある。
                全てのパターンは別々に計算されるから、グラフ上で5をコピーして2つにしても計算結果は変わらない。
                ここでは新しく10,11ノードを作っている。
                \par
                同じことを1ノードに対しても行っている。ここでは新しく12,13ノードを作っている。\\
                \noindent
                \textbf{Step4}
                \par
                $s_7$からは必ず10を通るから、グラフ上で$s_7\rightarrow 10$を予め計算しておいても問題ない。
                $s_5$についても同じことをしている。
                \par
                図4には3つのループがあるが、ループ上に割り込みが無いので、グラフ上で1本の辺に単純化しても計算結果は変わらない。\\
                \noindent
                \textbf{Step5}
                \par
                $s_6$と$s_7$を統合して$s_8$を作った。Step1でやったことと同じ。\\
                \quad\par
                ここまでくると、次のことが証明できれば定理が証明できたことになるのがわかる。\\
                \begin{shadebox}
                    有向グラフ$G(V,E)$を考える。ある頂点$s$と$g$を選ぶ。これらは次の条件を満たすものとする。
                    \begin{itemize}
                        \item $s$は$g$に隣接する。
                        \item $g$を含む閉路が存在したとしても、その閉路上に$s$は存在しない
                    \end{itemize}
                    このとき、$S$の点から$g$へ至る経路全ての連鎖重みを足し上げると、結果は1以下である。
                    \begin{figure}[H]
                        \centering
                        \includegraphics[keepaspectratio, scale=0.7]
                        {\currfiledir/imgs/LemmaForSumOfChainedWeightPath/img.eps}
                        \caption{$G(V,E)$}
                    \end{figure}
                \end{shadebox}
                この補題の証明は次の小小節で行う。ここではこの補題を認めることにする。
                この補題を図6の2ノードに対して適用すると、$s_5$ノードから2ノードへの全てのパスについて連鎖的重みを足し合わせると1になることがわかる。
                よって2ノードを新しく、例えば$s_9$で置き換えてから、さらに$s_8$と統合させて再び補題を適用すると定理が証明できたことになる。
                \par
                上の図の場合は連鎖的重みの全ての経路についての総和がちょうど1になったが、これは$S$の点から$g$への全ての経路上のノードにおいて、「$S$からも$g$からも入れないような経路」からの割り込みがなかったからである。
                もし、例えば8に$x$ノードが隣接し、$x$の入次数が0であると、辺$x\rightarrow8$の重みのせいで他の辺の重みが下がる(重み和は1でないといけなかった！)から結果は1を下回る。
            \end{proof}
        \subsection{(補題) ゴールに隣接する1つのスタート地点からの全ての経路の連鎖重みの総和は1以下}
        \begin{shadebox}
            有向グラフ$G(V,E)$を考える。ある頂点$s$と$g$を選ぶ。これらは次の条件を満たすものとする。
            \begin{itemize}
                \item $s$は$g$に隣接する。
                \item $g$を含む閉路が存在したとしても、その閉路上に$s$は存在しない
            \end{itemize}
            このとき、$S$の点から$g$へ至る経路全ての連鎖重みを足し上げると、結果は1以下である。
        \end{shadebox}
        状況を図に表すと次のようになる。
        \begin{figure}[H]
            \centering
            \includegraphics[keepaspectratio, scale=0.7]
            {\currfiledir/imgs/LemmaForSumOfChainedWeightPath/img.eps}
            \caption{$G(V,E)$}
        \end{figure}
        ゴール$g$に隣接するスタート$s$があり、$g$を含む閉路が$l$本存在する。
        もちろん$w_s + \sum_{i=1}^{l}{w_i} \leq 1$である。
        \begin{proof}
            \quad\par
            まず$s$から$g$へ来てそれでお終いというパターンの重みは$w_s$である。
            次に、一旦$g$へ来てから閉路を何度か回るパターンを全列挙するには、
            $l$本の閉路から重複を許して$k$本選んで回るパターンを$k=1,\dotsc,\infty$まで考えればよい。
            各パターンの重みは、最初に$s\rightarrow g$を通っていることに注意すると
            \begin{align*}
                w_s\sum_{k=1}^\infty{(w_1+\dotsb+w_l)^k} &= w_s\left(\frac{1}{1-(w_1+\dotsb+w_l)}-1\right) = w_s\frac{w_1+\dotsb+w_l}{1-(w_1+\dotsb+w_l)}\\
                &\leq w_s\frac{w_1+\dotsb+w_l}{w_s}\quad(\because\;w_s+w_1+\dotsb+w_l\leq 1)\\
                &= w_1+\dotsb+w_l
            \end{align*}
            よって、最初に考えた$s\rightarrow g$で終わりのパターンと合わせると総和は$w_s+w_1+\dotsb+w_l\leq 1$となる。
        \end{proof}